\chapter{基本概念}
\orig{50--68}

\section{映射、变换}

设 $A_1,A_2,\ldots,A_n$ 是 $n$ 个集合，$D$ 是另一个集合。

\textbf{定义。} 若通过一个法则 $\phi$，对 $A_1\times A_2\times\cdots\times A_n$ 的每一个元
\[
(a_1,a_2,\ldots,a_n),\qquad a_i\in A_i,
\]
都得到唯一的一个 $D$ 中的元 $d$，那么这个法则 $\phi$ 叫做集合 $A_1\times A_2\times\cdots\times A_n$ 到集合 $D$ 的一个\textbf{映射}，记作
\[
\phi:(a_1,a_2,\ldots,a_n)\longmapsto d=\phi(a_1,a_2,\ldots,a_n).
\]

\textbf{定义。} $A\times B$ 到 $D$ 的一个映射叫做 $A\times B$ 到 $D$ 的一个\textbf{代数运算}。一个代数运算也可写成
\[
(a,b)\longmapsto d=a\circ b.
\]

\textbf{定义。} 在集合 $A$ 到集合 $\bar A$ 的一个映射 $\phi$ 下，$\bar A$ 的每一个元都至少是 $A$ 中某一个元的像，那么 $\phi$ 叫做 $A$ 到 $\bar A$ 的一个\textbf{满射}。

\textbf{定义。} $A$ 到 $\bar A$ 的一个映射
\[
\phi:a\longmapsto \bar a
\]
叫做 $A$ 到 $\bar A$ 的一个\textbf{单射}，如果
\[
a\ne b\Longrightarrow \bar a\ne\bar b.
\]

\textbf{定义。} 若 $A$ 到 $\bar A$ 的一个映射既是满射又是单射，则称 $\phi$ 为 $A$ 与 $\bar A$ 间的一个\textbf{一一映射}。

\textbf{定义。} $A$ 到 $A$ 的映射叫做 $A$ 的一个\textbf{变换}。$A$ 到 $A$ 的满射、单射或一一映射，分别叫做 $A$ 的一个\textbf{满射变换}、\textbf{单射变换}或\textbf{一一变换}。

下面给出若干反例。

\textbf{1.} 令 $A=D=\mathbb R$，写
\[
\phi:\quad
\begin{cases}
a\longmapsto a,&a\ne1,\\
1\longmapsto b,&b^2=1.
\end{cases}
\]
这不是 $A$ 到 $D$ 的一个映射，因为 $1$ 的像不是唯一的。

\textbf{2.} 令 $A=D=\mathbb R$，写
\[
\phi:a\longmapsto\sqrt a.
\]
这也不是 $A$ 到 $D$ 的一个映射，因为当 $a<0$ 时，$\sqrt a\notin D$。

\textbf{3.} 令
\[
A=\{1,2,3\},\qquad D=\{1,2,\ldots,16\},
\]
并定义
\[
\phi_1(a)=2^a,\qquad \phi_2(a)=a^2-a+2.
\]
容易验证 $\phi_1$ 与 $\phi_2$ 是 $A$ 到 $D$ 的两个相同的映射。

\textbf{4.} 若改取
\[
A=\{1,2,3,4\},\qquad D=\{1,2,\ldots,16\},
\]
仍令
\[
\phi_1(a)=2^a,\qquad \phi_2(a)=a^2-a+2,
\]
则 $\phi_1$ 与 $\phi_2$ 都是 $A$ 到 $D$ 的映射，但不相等，因为
\[
\phi_1(4)=2^4\ne4^2-4+2=\phi_2(4).
\]

\textbf{5.} 令
\[
A=\{1,2,3\},\qquad \bar A=\{4,5,6,7\},
\]
定义
\[
1\mapsto4,
\qquad2\mapsto5,
\qquad3\mapsto6.
\]
这是单射而不是满射，故不是一一映射。

\textbf{6.} 令
\[
A=\{1,2,3\},\qquad \bar A=\{0\},
\]
定义 $a\mapsto0$。这是满射而不是单射，故不是一一映射。

\textbf{7.} 令 $A=\mathbb R$，
\[
\tau:x\longmapsto e^x.
\]
这是 $A$ 的一个单射变换，不是满射变换，故不是一一变换。

\textbf{8.} 令 $A=\mathbb Z$，定义
\[
\tau(a)=
\begin{cases}
\dfrac a2,&a\text{ 为偶数},\\[4pt]
\dfrac{a+1}{2},&a\text{ 为奇数}.
\end{cases}
\]
这是 $A$ 的一个满射变换，不是单射变换，故不是一一变换。

\textbf{9.} 令 $A=\mathbb Z$，$\bar A=2\mathbb Z$，定义
\[
\phi:n\longmapsto2n.
\]
$\phi$ 是 $A$ 与 $\bar A$ 间的一一映射，而 $\bar A$ 是 $A$ 的真子集。因此不能说 $\phi$ 是 $A$ 的一一变换；若把同一法则看作 $A\to A$ 的映射，则它只是 $A$ 的单射变换。

\textbf{10.} 令 $\bar A=2\mathbb Z$，$A=\mathbb Z$，定义
\[
\phi:2n\longmapsto n.
\]
这是 $\bar A$ 与 $A$ 间的一一映射，而 $\bar A$ 是 $A$ 的真子集，所以它不是 $A$ 的变换，更不是 $A$ 的一一变换。

\textbf{11.} 令
\[
\omega=-\frac12+\frac{\sqrt3}{2}i,
\]
并置
\[
A=\{a+b\omega+c\omega^2\mid a,b,c\in\mathbb Q\},
\qquad
\bar A=\{x+y\sqrt2\mid x,y\in\mathbb Q\}.
\]
定义
\[
\phi(a+b\omega+c\omega^2)=(2a-b-c)+(b-c)\sqrt2.
\]
下面说明 $\phi$ 是 $A$ 到 $\bar A$ 的一一映射。

首先验证它确实是映射。若
\[
a+b\omega+c\omega^2=d+e\omega+f\omega^2,
\]
则利用
\[
\omega=-\frac12+\frac{\sqrt3}{2}i,
\qquad
\omega^2=-\frac12-\frac{\sqrt3}{2}i,
\]
比较实部与虚部，得到
\[
a-\frac12(b+c)=d-\frac12(e+f),
\qquad
b-c=e-f.
\]
也就是
\[
2a-b-c=2d-e-f,
\qquad
b-c=e-f,
\]
所以两种表示在 $\phi$ 下有相同的像。这一点是必要的，因为 $1+\omega+\omega^2=0$，同一元素的三元系数表示并不唯一。

其次验证满射。任取
\[
p+q\sqrt2\in\bar A,
\]
有
\[
\frac{p+q}{2}+q\omega+0\omega^2
\longmapsto p+q\sqrt2.
\]

最后验证单射。若
\[
(2a-b-c)+(b-c)\sqrt2=(2a'-b'-c')+(b'-c')\sqrt2,
\]
则
\[
2a-b-c=2a'-b'-c',
\qquad
b-c=b'-c'.
\]
由此得到
\[
a-c=a'-c',
\qquad
 a-b=a'-b'.
\]
于是
\begin{align*}
a+b\omega+c\omega^2
&=a+(a-a'+b')\omega+(a-a'+c')\omega^2\\
&=(a+a\omega+a\omega^2)-a'(\omega+\omega^2)+b'\omega+c'\omega^2\\
&=a'+b'\omega+c'\omega^2,
\end{align*}
其中使用了 $1+\omega+\omega^2=0$。故 $\phi$ 是单射，从而是一一映射。

\section{基本算律}

\textbf{定义。} 一个集合 $A$ 的代数运算 $\circ$ 称为适合\textbf{结合律}，若对任意 $a,b,c\in A$ 都有
\[
(a\circ b)\circ c=a\circ(b\circ c).
\]

\textbf{定义。} 一个 $A\times A$ 到 $D$ 的代数运算称为适合\textbf{交换律}，若对任意 $a,b\in A$ 都有
\[
a\circ b=b\circ a.
\]

以下给分配律以定义。设 $\odot$ 是一个 $B\times A$ 到 $A$ 的代数运算，$\oplus$ 是 $A$ 的代数运算。若
\[
b\odot(a_1\oplus a_2)=(b\odot a_1)\oplus(b\odot a_2),
\]
则称 $\odot$ 对 $\oplus$ 适合\textbf{第一（左）分配律}。第二个（右）分配律的定义完全类似。

\subsection*{关于结合律}

\textbf{1.} $A=\mathbb Z$，运算为普通减法。它不适合结合律，因为当 $c\ne0$ 时
\[
(a-b)-c\ne a-(b-c).
\]

\textbf{2.} $A=\mathbb R\setminus\{0\}$，令
\[
a\circ b=\frac ab.
\]
该运算不适合结合律，例如
\[
(1\circ1)\circ2=\frac12,
\qquad
1\circ(1\circ2)=2.
\]

\textbf{3.} $A=\mathbb R$，令
\[
a\circ b=a+2b.
\]
则
\[
(a\circ b)\circ c=a+2b+2c,
\qquad
 a\circ(b\circ c)=a+2b+4c,
\]
所以当 $c\ne0$ 时结合律不成立。

\textbf{4.} $A=\mathbb Z_{>0}$，令
\[
a\circ b=a^b.
\]
该运算不适合结合律，因为
\[
(2\circ1)\circ3=(2^1)^3=8,
\qquad
2\circ(1\circ3)=2^{1^3}=2.
\]

\textbf{5.} 一个集合 $A$ 的代数运算若适合结合律，则对 $A$ 中任意 $n\ge3$ 个元 $a_1,a_2,\ldots,a_n$，用各种加括号的步骤所得结果都一样；但反之不真。

例如取 $A=\{a,b,c,d\}$，代数运算由表
\[
\begin{array}{c|cccc}
\circ&a&b&c&d\\\hline
 a&b&c&d&d\\
 b&d&d&d&d\\
 c&d&d&d&d\\
 d&d&d&d&d
\end{array}
\]
给出。任意四个元的任意加括号乘积都为 $d$，即“四元结合”成立；但三元结合律不成立，因为
\[
a\circ(a\circ a)=a\circ b=c,
\qquad
(a\circ a)\circ a=b\circ a=d.
\]

\textbf{6.} 下面给出一个对任意 $n>3$，“$n$ 元结合律”均不成立的例。取 $A=\{a,b,c\}$，运算表为
\[
\begin{array}{c|ccc}
\circ&a&b&c\\\hline
 a&b&c&a\\
 b&c&c&c\\
 c&c&c&c
\end{array}
\]
三元结合律已经不成立：
\[
(a\circ b)\circ b=c\circ b=c,
\qquad
 a\circ(b\circ b)=a\circ c=a.
\]
对任意 $n>3$，同样比较
\[
(a\circ b)\circ(b\circ\cdots\circ b)=c,
\qquad
 a\circ\bigl(b\circ(b\circ\cdots\circ b)\bigr)=a,
\]
可见相应的多元结合律也不成立。

\subsection*{关于交换律}

\textbf{1.} $A=\mathbb R$，普通减法 $a\circ b=a-b$ 不适合交换律；当 $a\ne b$ 时一般有
\[
a-b\ne b-a.
\]

\textbf{2.} 取 $A=\{a,b,c,d\}$，代数运算由
\[
\begin{array}{c|cccc}
\circ&a&b&c&d\\\hline
 a&a&b&c&d\\
 b&b&d&a&c\\
 c&c&a&b&d\\
 d&d&c&a&b
\end{array}
\]
给出。该运算不适合交换律，例如
\[
c\circ d=d,
\qquad
 d\circ c=a.
\]

\textbf{3.} 若集合 $A$ 的代数运算同时适合结合律与交换律，那么计算
\[
a_1\circ a_2\circ\cdots\circ a_n
\]
时可以任意掉换元素次序；但反之不真。

例如仍取 $A=\{a,b,c,d\}$，令
\[
\begin{array}{c|cccc}
\circ&a&b&c&d\\\hline
 a&d&c&d&d\\
 b&d&d&d&d\\
 c&d&d&d&d\\
 d&d&d&d&d
\end{array}
\]
任意三个元之积均为 $d$，因而三元结合律及三元交换律成立；但二元交换律不成立，因为
\[
a\circ b\ne b\circ a.
\]

\textbf{4.} 取 $A=\{a,b,c\}$，运算表为
\[
\begin{array}{c|ccc}
\circ&a&b&c\\\hline
 a&b&c&c\\
 b&c&b&b\\
 c&c&b&b
\end{array}
\]
从表中可见二元交换律成立，但三元结合律不成立：
\[
(a\circ b)\circ c=c\circ c=b,
\qquad
 a\circ(b\circ c)=a\circ b=c.
\]
三元交换律也不成立，例如
\[
(a\circ b)\circ c=b,
\qquad
 a\circ(c\circ b)=c.
\]

\textbf{5.} 三元结合律成立、二元交换律不成立，并不能推出多元结合交换律。例如全体 $n$ 阶方阵在矩阵乘法下满足结合律，但一般不满足交换律；若 $AB\ne BA$，则
\[
(EA)B\ne E(BA).
\]

\textbf{6.} 取 $A=\{a,b\}$，运算表
\[
\begin{array}{c|cc}
\circ&a&b\\\hline
 a&a&a\\
 b&b&a
\end{array}
\]
满足三元结合律，但多元交换律不成立。实际上
\[
a\circ b\circ a\circ\cdots\circ a=a,
\qquad
 b\circ a\circ a\circ\cdots\circ a=b,
\]
故改变因子的次序可能改变结果。

\subsection*{关于分配律}

若 $\oplus$ 适合结合律，且 $\odot$ 对 $\oplus$ 适合第一（左）分配律，那么
\[
b\odot(a_1\oplus\cdots\oplus a_n)
=(b\odot a_1)\oplus\cdots\oplus(b\odot a_n),
\]
但反之不真。

例如取 $A=\{a,b,c\}$，令 $\oplus$ 与 $\odot$ 分别由
\[
\begin{array}{c|ccc}
\oplus&a&b&c\\\hline
 a&b&c&c\\
 b&c&c&c\\
 c&c&c&c
\end{array}
\qquad
\begin{array}{c|ccc}
\odot&a&b&c\\\hline
 a&a&c&c\\
 b&c&c&c\\
 c&c&c&c
\end{array}
\]
给出。容易验证 $\oplus$ 适合三元结合律，并且 $\odot$ 对 $\oplus$ 满足三元分配式
\[
b\odot(a_1\oplus a_2\oplus a_3)
=(b\odot a_1)\oplus(b\odot a_2)\oplus(b\odot a_3),
\]
但二元分配律并不成立，因为
\[
a\odot(a\oplus a)=a\odot b=c,
\]
而
\[
(a\odot a)\oplus(a\odot a)=a\oplus a=b.
\]

\section{同态、同构}

\textbf{定义。} 一个 $A$ 到 $\bar A$ 的映射 $\phi$，若对代数运算 $\circ$ 和 $\bar\circ$ 来说满足：只要
\[
a\longmapsto\bar a,
\qquad b\longmapsto\bar b,
\]
就有
\[
a\circ b\longmapsto\bar a\,\bar\circ\,\bar b,
\]
则称 $\phi$ 为关于这两个代数运算的\textbf{同态映射}。

\textbf{定义。} 若关于 $\circ$ 和 $\bar\circ$ 有一个 $A$ 到 $\bar A$ 的满射同态映射，则称它为\textbf{同态满射}，并说 $A$ 与 $\bar A$ 关于这两个运算\textbf{同态}。

\textbf{定义。} 若 $A$ 与 $\bar A$ 间的一一映射 $\phi$ 同时是关于 $\circ$ 和 $\bar\circ$ 的同态映射，则称 $\phi$ 为\textbf{同构映射}，并说 $A$ 与 $\bar A$ 同构，记为
\[
A\cong\bar A.
\]

\textbf{定义。} 关于 $\circ$ 的一个 $A$ 与 $A$ 间的同构映射叫做 $A$ 的一个\textbf{自同构}。

\textbf{1.} 令 $A=\mathbb Z$，$A$ 的代数运算是普通加法；令 $\bar A=\{1,-1\}$，$\bar A$ 的代数运算是普通乘法。映射
\[
\phi:a\longmapsto-1
\]
虽是 $A$ 到 $\bar A$ 的映射，却不是同态映射，因为
\[
a\mapsto-1,
\quad b\mapsto-1,
\quad a+b\mapsto-1\ne(-1)(-1).
\]

\textbf{2.} 仍令 $A=\mathbb Z$（普通加法），$\bar A=\{1,-1\}$（普通乘法），定义
\[
\phi(a)=
\begin{cases}
1,&a\text{ 为偶数},\\
-1,&a\text{ 为奇数}.
\end{cases}
\]
这是 $A$ 到 $\bar A$ 的同态满射，但不是同构映射。

\textbf{3.} 若改为
\[
\phi(a)=
\begin{cases}
-1,&a\text{ 为偶数},\\
1,&a\text{ 为奇数},
\end{cases}
\]
则它不是同态满射。例如
\[
2+2=4\mapsto-1\ne(-1)(-1).
\]

这说明：即使某个映射不是同态满射，也不妨碍存在另一个同态满射；只要存在一个满射同态映射，就说相应的两个代数系统同态。

\textbf{4.} 再取 $A=\mathbb Z$，但这次 $A$ 的代数运算为普通乘法；$\bar A=\{1,-1\}$ 仍用普通乘法。定义
\[
\psi(a)=
\begin{cases}
1,&a\text{ 为偶数},\\
-1,&a\text{ 为奇数}.
\end{cases}
\]
虽然 $\psi$ 是满射，却不是同态映射，因为
\[
2\mapsto1,
\qquad 3\mapsto-1,
\qquad 2\cdot3=6\mapsto1\ne1\cdot(-1).
\]
同一个集合之间的映射能否成为同态，不能脱离所指定的代数运算来判断。

\textbf{5.} 令 $A=\mathbb R$，$A$ 的代数运算是普通乘法。考察下列映射是否是 $A$ 到某个子集 $\bar A$ 的同态满射：
\[
\text{a) }x\mapsto|x|,
\qquad
\text{b) }x\mapsto2x,
\qquad
\text{c) }x\mapsto x^2,
\qquad
\text{d) }x\mapsto-x.
\]
答案为：

\textbf{a)} $x\mapsto|x|$ 是 $A$ 到 $\bar A=\mathbb R_{\ge0}$ 的同态满射，但不是同构映射。

\textbf{b)} 由于
\[
\phi(xy)=2xy\ne(2x)(2y)=\phi(x)\phi(y)
\]
（除非 $xy=0$），故 $x\mapsto2x$ 不是同态满射。

\textbf{c)} $x\mapsto x^2$ 是 $A$ 到 $\mathbb R_{\ge0}$ 的同态满射，但不是同构映射。

\textbf{d)} 一般有
\[
-(xy)\ne(-x)(-y),
\]
故 $x\mapsto-x$ 不是同态满射。

这个例子说明：同态满射既要是满射，又要保持运算；同构还要求是一一映射。

\textbf{6.} 令 $A=\{a,b,c\}$，代数运算由
\[
\begin{array}{c|ccc}
\circ&a&b&c\\\hline
 a&c&c&c\\
 b&c&c&c\\
 c&c&c&c
\end{array}
\]
给出。$A$ 的一一变换共有六个：
\begin{align*}
\tau_1&:a\mapsto a,
&b\mapsto b,
&c\mapsto c;\\
\tau_2&:a\mapsto b,
&b\mapsto a,
&c\mapsto c;\\
\tau_3&:a\mapsto b,
&b\mapsto c,
&c\mapsto a;\\
\tau_4&:a\mapsto c,
&b\mapsto b,
&c\mapsto a;\\
\tau_5&:a\mapsto c,
&b\mapsto a,
&c\mapsto b;\\
\tau_6&:a\mapsto a,
&b\mapsto c,
&c\mapsto b.
\end{align*}
容易验证 $\tau_1,\tau_2$ 是 $A$ 的自同构，而 $\tau_3,\tau_4,\tau_5,\tau_6$ 不是。

\textbf{7.} 假定关于代数运算 $\circ$ 与 $\bar\circ$，$A$ 与 $\bar A$ 同态，那么：

(i) 若 $\circ$ 适合结合律，则 $\bar\circ$ 也适合结合律；

(ii) 若 $\circ$ 适合交换律，则 $\bar\circ$ 也适合交换律；但两者的逆命题都不真。

\textbf{例 1。} 令 $A=\mathbb Q$，定义
\[
a\circ b=-a-b;
\]
令 $\bar A=\{1\}$，$\bar\circ$ 为普通乘法，并令 $\phi(a)=1$。显然 $\phi$ 是同态满射，$\bar\circ$ 适合结合律；但 $\circ$ 不适合结合律，例如
\[
(2\circ1)\circ1=(-2-1)\circ1=3-1=2,
\]
而
\[
2\circ(1\circ1)=2\circ(-1-1)=-2-(-2)=0.
\]

\textbf{例 2。} 令 $A=\mathbb Q$，定义
\[
a\circ b=b;
\]
仍令 $\bar A=\{1\}$，$\bar\circ$ 为普通乘法，$\phi(a)=1$。此时 $\bar\circ$ 适合交换律，而 $\circ$ 不适合交换律，例如
\[
3\circ2=2\ne3=2\circ3.
\]

\textbf{8.} 设 $\odot,\oplus$ 都是集合 $A$ 的代数运算，$\bar\odot,\bar\oplus$ 都是集合 $\bar A$ 的代数运算，并存在一个 $A$ 到 $\bar A$ 的满射 $\phi$，使 $\phi$ 同时给出两对运算的同态。则：

(i) 若 $\odot$ 对 $\oplus$ 适合第一分配律，$\bar\odot$ 对 $\bar\oplus$ 也适合第一分配律；

(ii) 若 $\odot$ 对 $\oplus$ 适合第二分配律，$\bar\odot$ 对 $\bar\oplus$ 也适合第二分配律；反之不真。

例如令 $A=\mathbb Q$，定义
\[
a\odot b=a+b,
\qquad
a\oplus b=a+b,
\]
令 $\bar A=\{0\}$，$\bar\odot$ 为普通乘法，$\bar\oplus$ 为普通加法，并令 $\phi(a)=0$。对两对运算而言，$A$ 与 $\bar A$ 都同态。显然 $\bar\odot$ 对 $\bar\oplus$ 满足左右两个分配律，但 $\odot$ 对 $\oplus$ 两个分配律都不满足，例如
\[
1\odot(1\oplus1)=3\ne4=(1\odot1)\oplus(1\odot1),
\]
又
\[
(1\oplus1)\odot1=3\ne4=(1\odot1)\oplus(1\odot1).
\]

\textbf{9.} 令 $A=\mathbb Q$，$A$ 的运算为普通加法；令
\[
\bar A=\mathbb Q\setminus\{0\},
\]
$\bar A$ 的运算为普通乘法。可以证明，对这两个给定的代数运算来说，$A$ 与 $\bar A$ 之间不存在同构映射。

证明如下。假设 $\phi:A\to\bar A$ 是同构。加法单位元 $0$ 必须映到乘法单位元 $1$。由于 $\phi$ 是满射，存在 $x\in\mathbb Q$ 使
\[
\phi(x)=-1.
\]
又 $x\ne0$。于是
\[
\phi(2x)=\phi(x+x)=\phi(x)^2=(-1)^2=1=\phi(0).
\]
由于 $\phi$ 是单射，得到 $2x=0$，从而 $x=0$，矛盾。

也可更直接地观察：若 $\phi$ 是同构，则任意 $a\in\mathbb Q$ 都有
\[
\phi(a)=\phi\left(\frac a2+\frac a2\right)=\phi\left(\frac a2\right)^2>0,
\]
所以负有理数不可能出现在像中，与满射矛盾。

\section{等价关系}

设 $A$ 是一个集合，$D=\{\text{对},\text{错}\}$。

\textbf{定义。} $A\times A$ 到 $D$ 的一个映射 $R$ 叫做 $A$ 的元素间的一个\textbf{关系}。若 $R(a,b)=\text{对}$，称 $a$ 与 $b$ 符合关系 $R$，记作 $aRb$；若 $R(a,b)=\text{错}$，则称 $a$ 与 $b$ 不符合关系 $R$。

\textbf{定义。} 集合 $A$ 的元素间的一个关系 $\sim$ 叫做一个\textbf{等价关系}，若它满足：
\begin{enumerate}[label=\Roman*.,leftmargin=3em]
\item 反射律：$a\sim a$，对每个 $a\in A$；
\item 对称律：$a\sim b\Longrightarrow b\sim a$；
\item 推移律：$a\sim b,\ b\sim c\Longrightarrow a\sim c$。
\end{enumerate}
下面的例题说明这三个条件是相互独立的。

\textbf{1.} 有人试图由对称律和推移律推出反射律：由 $aRb$ 得 $bRa$，再由 $aRb,bRa$ 得 $aRa$。错误在于，这里的 $a$ 已经受到“存在 $b$ 使 $aRb$”这一条件的限制，并不是集合中任意的 $a$。

取 $A=\{a,b,c\}$，定义关系表
\[
\begin{array}{c|ccc}
R&a&b&c\\\hline
 a&\text{错}&\text{错}&\text{错}\\
 b&\text{错}&\text{对}&\text{对}\\
 c&\text{错}&\text{对}&\text{对}
\end{array}
\]
可直接验证 $R$ 满足对称律与推移律，但 $aRa$ 为错，所以不满足反射律。这说明反射律是独立的。

\textbf{2.} 令
\[
A=\{1,2,4,6,10\},
\qquad
aRb\Longleftrightarrow4\mid(a+b).
\]
$R$ 显然满足对称律。它也满足推移律：若 $aRb$ 与 $bRc$，则 $a\equiv c\equiv-b\pmod4$；在集合 $A$ 中，能与别的元素发生关系的余数类只有 $0$ 与 $2$，从而 $a+c\equiv0\pmod4$，即 $aRc$。更直接地说，关系的非空块为 $\{4\}$ 与 $\{2,6,10\}$。但
\[
4\nmid1+1,
\]
所以 $1$ 与自身不符合关系，$R$ 不满足反射律。
\jiaokan{此处用 $a=4q_1+\gamma_1$、$c=4q_3+\gamma_3$ 作分类讨论时，符号与余数限制的书写较为含混。结论本身正确；这里改用同余类直接验证推移律，并保留集合与关系不变。}

\textbf{3.} 令 $A=\{a,b\}$，定义
\[
\begin{array}{c|cc}
R&a&b\\\hline
 a&\text{对}&\text{错}\\
 b&\text{对}&\text{对}
\end{array}
\]
则 $R$ 满足反射律，也满足推移律，但不满足对称律。这说明对称律是独立的。

\textbf{4.} 令 $A=\mathbb R$，关系 $\ge$ 满足反射律和推移律，但不满足对称律。

\textbf{5.} 令 $A=\{a,b,c\}$，定义
\[
\begin{array}{c|ccc}
R&a&b&c\\\hline
 a&\text{对}&\text{对}&\text{对}\\
 b&\text{对}&\text{对}&\text{错}\\
 c&\text{对}&\text{错}&\text{对}
\end{array}
\]
则 $R$ 满足反射律和对称律，但不满足推移律：有 $bRa$ 且 $aRc$，却没有 $bRc$。这说明推移律也是独立的。
